\scanpage{40}
\[
 gh^2g^{-1}=\deleted{\unsure{u}}(\lambda h)^2=\lambda^2h^2=h^2
\]
donc $g\in C(h^2)$ donc o.k. si $h^2\ne1$ i.e. $h$ non une réflexion

Mais ceci nous fait une belle jambe dans le cas (typique de car 2 !) où \deleted{\uline{\unsure{tte section}}} $H$ lui-$=$ est tel que $h^2=1\quad\forall h\in\Gamma(U,H)$\ldots!

\deleted{\unsure{Il}} Mais alors est-il possible d'avoir $gug^{-1}=\lambda_u u\deleted{\unsure{{}^{\prime}|\cdot|{}^{-1}}}$ \deleted{\(\forall\)} pour $g$ \uline{fixé}, et tt $u$ \inserted{inv.} de la \deleted{\unsure{le fibré}} forme $\alpha+\beta u_0$

Prenant la section ``universelle'' $\alpha+\beta u_0$, on doit avoir \unsure{(identiquemt|identiquement} en $\alpha,\beta$
\[
 \underbrace{g(\alpha+\beta u_0)g^{-1}}_{\begin{gathered}
 \alpha+\beta gu_0g^{-1}=\\[-2pt]
 =\alpha+\beta\lambda_0u_0
 \end{gathered}}
 =\lambda(\alpha,\beta)(\alpha+\beta u_0)
\]
ce qui implique, \deleted{\unsure{faisant passer}}
\[
 \deleted{gu_0g^{-1}=\lambda_0u_0}
\]
\[
\begin{aligned}
 \alpha\lambda(\alpha,\beta)&\equiv1 \quad\text{(identité en $\alpha,\beta$ dans}\\
 &\mathcal O_S[\alpha,\beta]_{\Delta(\alpha,\beta)}
 \quad\underset{\text{localisation}}{\curvearrowright}
\end{aligned}
\]
d'où $\lambda(\alpha,\beta)=1$, o.k.

Donc $H$ est \uline{son propre centralisateur} (et aussi le centralisateur de tte section partout non de carré 1)

Est-il le centralisateur de $\mathfrak h$ ?
% archive sheet number: 39
